\documentclass[11pt]{article}

\usepackage{amsmath, amssymb, amsthm}
\usepackage{stmaryrd}   % \llbracket, \rrbracket, \fatsemi
\usepackage{mathtools}
\usepackage{booktabs}
\usepackage{graphicx}
\usepackage[margin=1in]{geometry}
\usepackage{hyperref}

% --- notation ---------------------------------------------------------------
\newcommand{\sem}[1]{\llbracket #1 \rrbracket}            % denotation
\newcommand{\Empty}{\bot}                                 % empty / undefined
\newcommand{\size}[1]{\lvert #1 \rvert}
\newcommand{\R}[5]{\langle #1,\, #2,\, #3,\, #4,\, #5 \rangle} % split repr.
\newcommand{\gmin}{g_{\min}}
\newcommand{\gmax}{g_{\max}}

% operators
\newcommand{\Then}{\mathbin{\fatsemi}}                    % THEN
\newcommand{\ThenG}[2]{\mathbin{\fatsemi_{[#1,\,#2]}}}    % THEN with gap
\newcommand{\opname}[1]{\mathbin{\textsc{#1}}}
\newcommand{\Ovl}{\opname{ovl}}                           % OVERLAP
\newcommand{\Cont}{\opname{cont}}                         % CONTAINS
\newcommand{\Eq}{\opname{eq}}                             % EQUAL

\theoremstyle{definition}
\newtheorem{definition}{Definition}
\theoremstyle{plain}
\newtheorem{theorem}{Theorem}

\title{Interval-Set Operators on the Split Representation\\
       \large definitions and machine-checked correctness}
\author{}
\date{}

\begin{document}
\maketitle

\noindent
This note collects, in human-readable form, the binary operators defined on the
compact \emph{split interval set} representation, together with the equivalences
proved in Lean~4 / Mathlib (one file per operator under
\texttt{lean-prove/Intervalset/}).  Every theorem is fully formal and uses only
the standard axioms \texttt{propext}, \texttt{Classical.choice},
\texttt{Quot.sound}.

% ===========================================================================
\section{Preliminaries}

\begin{definition}[Interval]
An \emph{interval} $I = [a, b]$ is a pair of naturals with $a \le b$; its
\emph{size} is $\size{I} = b - a + 1 \ge 1$.
\end{definition}

\begin{definition}[Split interval set]
A \emph{split interval set} $R = \R{s}{l}{r}{e}{c}$ denotes all intervals that
start in $[s, l]$, end in $[r, e]$, and have size at least $c$:
\[
  \sem{R} = \bigl\{\, [a, b] \;\bigm|\;
            s \le a \le l,\;\; r \le b \le e,\;\; b - a + 1 \ge c \,\bigr\}.
\]
In the $(\text{start},\text{stop})$ plane this is the box $[s,l]\times[r,e]$ with
the lower-left corner clipped by the diagonal $\text{stop}-\text{start}+1 \ge c$.
\end{definition}

\paragraph{Enumeration semantics.}
Each operator is first given on \emph{single} intervals as a partial map
$I_1 \mathbin{\textsc{op}} I_2 \in \{[\,\cdot\,,\cdot\,], \Empty\}$.  It is lifted
to sets by enumerating all pairs and keeping the non-empty results:
\[
  A \mathbin{\textsc{op}} B \;=\;
  \bigl\{\, I_1 \mathbin{\textsc{op}} I_2 \;\bigm|\; I_1 \in A,\; I_2 \in B,\;
            I_1 \mathbin{\textsc{op}} I_2 \neq \Empty \,\bigr\}.
\]

\paragraph{Correctness statement.}
For each operator we define a function on split interval sets returning an
\texttt{Option} (a representation, or \texttt{none} for $\Empty$), and prove that,
for every interval $i$,
\[
  i \in \sem{A \mathbin{\textsc{op}} B}
  \quad\Longleftrightarrow\quad
  \exists\, I_1 \in \sem{A},\, I_2 \in \sem{B}.\;\;
  I_1 \mathbin{\textsc{op}} I_2 = i .
\]
Equivalently $\sem{A \mathbin{\textsc{op}} B} = \sem{A} \mathbin{\textsc{op}}
\sem{B}$: the compact operator computes exactly the enumeration semantics.

\paragraph{Conventions.}
We assume the relevant size bounds are $\ge 1$ where stated (harmless: every
interval has size $\ge 1$, so $c=0$ and $c=1$ denote the same set).  All
``$-$'' below are natural-number subtractions; the feasibility guards are chosen
so no subtraction is silently floored to $0$.  We write $A=\R{s_1}{l_1}{r_1}{e_1}{c_1}$
and $B=\R{s_2}{l_2}{r_2}{e_2}{c_2}$, and a result $\R{s'}{l'}{r'}{e'}{c'}$.

% ===========================================================================
\section{\textsc{Then} — concatenation (\emph{meets})}
\label{sec:then}
\textit{File: \texttt{Then.lean}, theorem \texttt{thenOp\_spec}; assumes
$c_1,c_2\ge1$.}

\[
  [a_1,b_1] \Then [a_2,b_2] =
  \begin{cases} [a_1,b_2] & a_2 = b_1+1,\\ \Empty & \text{otherwise.}\end{cases}
\]
The second interval must start on the very next frame (Allen's \emph{meets}).
When the feasibility guard
\[
  r_1 \le e_1,\;\; r_1+1 \le l_2,\;\; s_2 \le e_1+1,\;\; s_2 \le l_2,\;\; c_1 \le l_2
\]
holds,
\[
  A \Then B = \R{s_1}{\,l'\,}{\,r'\,}{e_2}{\,c_1+c_2\,},\quad
  \begin{aligned}
    l' &= \min(l_1,\; e_1+1-c_1,\; l_2-c_1),\\
    r' &= \max(r_2,\; r_1+c_2,\; s_2+c_2-1).
  \end{aligned}
\]
The result spans both intervals; its size is $c_1+c_2$.

% ===========================================================================
\section{\textsc{Then}$_{[\gmin,\gmax]}$ — bounded-gap concatenation (\emph{before})}
\label{sec:theng}
\textit{File: \texttt{ThenWithGap.lean}, theorem \texttt{thenGapOp\_spec}; assumes
$c_1,c_2\ge1$ and $\gmin\le\gmax$.}

\[
  [a_1,b_1] \ThenG{\gmin}{\gmax} [a_2,b_2] =
  \begin{cases} [a_1,b_2] & b_1+1+\gmin \le a_2 \le b_1+1+\gmax,\\
                \Empty & \text{otherwise,}\end{cases}
\]
i.e.\ a gap of $g\in[\gmin,\gmax]$ frames is allowed between the two intervals
(Allen's \emph{before} for $\gmin\ge1$; \textsc{Then} is the case
$\gmin=\gmax=0$, proved as \texttt{thenGapOp\_zero\_zero}).  When
\[
  r_1 \le e_1,\;\; r_1+\gmin+1 \le l_2,\;\; s_2 \le e_1+\gmax+1,\;\; s_2 \le l_2,\;\; c_1+\gmin \le l_2,
\]
\[
  A \ThenG{\gmin}{\gmax} B = \R{s_1}{\,l'\,}{\,r'\,}{e_2}{\,c_1+c_2+\gmin\,},\quad
  \begin{aligned}
    l' &= \min(l_1,\; e_1+1-c_1,\; l_2-c_1-\gmin),\\
    r' &= \max(r_2,\; r_1+c_2+\gmin,\; s_2+c_2-1).
  \end{aligned}
\]
The minimum gap $\gmin$ enters the size bound ($c'=c_1+c_2+\gmin$) and the
$B$-facing bounds; the start bound $e_1+1-c_1$ is \emph{not} relaxed by the gap,
since the first stop is still capped by $e_1$.

% ===========================================================================
\section{\textsc{Overlap} (Allen's \emph{overlaps})}
\label{sec:overlap}
\textit{File: \texttt{Overlap.lean}, theorem \texttt{overlapOp\_spec}; assumes
$c_1,c_2\ge1$.}

\[
  [a_1,b_1] \Ovl [a_2,b_2] =
  \begin{cases} [a_1,b_2] & a_1 < a_2 \le b_1 < b_2,\\
                \Empty & \text{otherwise.}\end{cases}
\]
The intervals properly overlap (second starts strictly later, ends strictly
later) and are merged.  When
\[
  r_1 \le e_1,\;\; s_2 \le e_1,\;\; s_2 \le l_2,\;\; s_1+1 \le e_1,\;\; s_1+1 \le l_2,
\]
\[
  A \Ovl B = \R{s_1}{\,l'\,}{\,r'\,}{e_2}{\,c'\,},\quad
  \begin{aligned}
    l' &= \min(l_1,\; e_1+1-c_1,\; e_1-1,\; l_2-1),\\
    r' &= \max(r_2,\; r_1+1,\; s_2+c_2-1,\; s_2+1),\\
    c' &= \max(c_1+1,\; c_2+1,\; 3).
  \end{aligned}
\]
The size floor $3$ is the smallest proper overlap $[a,a+2]$; the strictness
$a_1<a_2$ contributes the $e_1-1$, $l_2-1$ start caps and the $c_2+1$ size term.

% ===========================================================================
\section{\textsc{Contains} (Allen's \emph{contains})}
\label{sec:contains}
\textit{File: \texttt{Contains.lean}, theorem \texttt{containsOp\_spec}; assumes
$c_2\ge1$.}

\[
  [a_1,b_1] \Cont [a_2,b_2] =
  \begin{cases} [a_1,b_1] & a_1 < a_2 \;\wedge\; b_2 < b_1,\\
                \Empty & \text{otherwise.}\end{cases}
\]
The second interval lies strictly inside the first, which is kept.  When
\[
  s_2 \le l_2,\;\; r_2 \le e_2,\;\; s_1+1 \le l_2,\;\; c_2 \le e_2,
\]
\[
  A \Cont B = \R{s_1}{\,l'\,}{\,r'\,}{e_1}{\,c'\,},\quad
  \begin{aligned}
    l' &= \min(l_1,\; l_2-1,\; e_2-c_2),\\
    r' &= \max(r_1,\; r_2+1,\; s_2+c_2),\\
    c' &= \max(c_1,\; c_2+2).
  \end{aligned}
\]
The result inherits $A$'s start floor $s_1$ and stop cap $e_1$, but tightens the
start cap and stop floor on both sides to leave strict room for the enclosed
interval; enclosing a size-$c_2$ interval needs size $\ge c_2+2$.

% ===========================================================================
\section{\textsc{Equal} (Allen's \emph{equals})}
\label{sec:equal}
\textit{File: \texttt{Equal.lean}, theorem \texttt{equalsOp\_spec}.}

\[
  [a_1,b_1] \Eq [a_2,b_2] =
  \begin{cases} [a_1,b_1] & a_1=a_2 \;\wedge\; b_1=b_2,\\
                \Empty & \text{otherwise.}\end{cases}
\]
An interval is kept iff it occurs identically in both operands, so
$\sem{A \Eq B} = \sem{A}\cap\sem{B}$.  This is the field-wise intersection:
\[
  A \Eq B = \R{\max(s_1,s_2)}{\min(l_1,l_2)}{\max(r_1,r_2)}{\min(e_1,e_2)}{\max(c_1,c_2)},
\]
defined when $\max(s_1,s_2)\le\min(e_1,e_2)$ and
$\max(c_1,c_2)\le\min(e_1,e_2)-\max(s_1,s_2)+1$.  (This operator coincides with
set intersection / \textsc{And}.)

% ===========================================================================
\section{Summary}

\renewcommand{\arraystretch}{1.5}
\begin{table}[h]
\centering
\resizebox{\textwidth}{!}{%
\begin{tabular}{@{}llllll@{}}
\toprule
Op & $s'$ & $l'$ & $r'$ & $e'$ & $c'$ \\
\midrule
\textsc{Then}
  & $s_1$ & $\min(l_1, e_1{+}1{-}c_1, l_2{-}c_1)$
  & $\max(r_2, r_1{+}c_2, s_2{+}c_2{-}1)$ & $e_2$ & $c_1{+}c_2$ \\
\textsc{Then}$_{[\gmin,\gmax]}$
  & $s_1$ & $\min(l_1, e_1{+}1{-}c_1, l_2{-}c_1{-}\gmin)$
  & $\max(r_2, r_1{+}c_2{+}\gmin, s_2{+}c_2{-}1)$ & $e_2$ & $c_1{+}c_2{+}\gmin$ \\
\textsc{Overlap}
  & $s_1$ & $\min(l_1, e_1{+}1{-}c_1, e_1{-}1, l_2{-}1)$
  & $\max(r_2, r_1{+}1, s_2{+}c_2{-}1, s_2{+}1)$ & $e_2$ & $\max(c_1{+}1, c_2{+}1, 3)$ \\
\textsc{Contains}
  & $s_1$ & $\min(l_1, l_2{-}1, e_2{-}c_2)$
  & $\max(r_1, r_2{+}1, s_2{+}c_2)$ & $e_1$ & $\max(c_1, c_2{+}2)$ \\
\textsc{Equal}
  & $\max(s_1,s_2)$ & $\min(l_1,l_2)$ & $\max(r_1,r_2)$ & $\min(e_1,e_2)$ & $\max(c_1,c_2)$ \\
\bottomrule
\end{tabular}%
}
\caption{Compact result fields (subject to each operator's feasibility guard and
the well-formedness conditions $s'\le e'$, $c'\le e'-s'+1$).}
\end{table}

\noindent
Note how the result inherits its \emph{start} side ($s',l'$) from whichever
operand supplies the result's left endpoint and its \emph{stop} side ($r',e'$)
from whichever supplies the right endpoint: \textsc{Then}/\textsc{Then}$_{[\cdot]}$
and \textsc{Overlap} take $s_1$ from $A$ and $e_2$ from $B$; \textsc{Contains}
keeps both from $A$; \textsc{Equal} intersects.

\end{document}
